% ================= METHOD =================
\section{Method}
\label{sec:method}

\subsection{Blur model}
A uniform motion blur is modelled as $b = l * k + n$, where $l$ is the sharp image, $n$ is
noise and $k$ is a line kernel of length $L$ pixels at angle $\theta$ (counter-clockwise from
the $x$ axis). We generate $k$ by area sampling of the line segment, so non-integer lengths
and arbitrary angles are exact up to discretisation and an axis-aligned kernel is an exact box.
In the frequency domain $B = L\cdot K + N$ with
\begin{equation}
  |K(\omega)| = \left|\operatorname{sinc}\!\left(L\,\langle\omega, u_\theta\rangle\right)\right|,
\end{equation}
which vanishes on parallel lines perpendicular to the motion direction $u_\theta$, spaced
$1/L$ cycles per pixel apart. Both blur parameters are encoded in this null pattern: its
orientation gives $\theta$ and its spacing gives $L$.

\subsection{Radon--cepstrum estimator}
The estimator works on the log-magnitude spectrum $S = \log(\epsilon + |\mathcal{F}\{w\cdot g\}|)$
of the windowed (Tukey), mean-removed grayscale image $g$, zero-padded to a square grid of at
least $256\times256$ so that all directions are sampled isotropically. The spectrum is median
centred and multiplied by a raised-cosine disc that ends at a radial band limit $c$ (a fraction
of Nyquist; $c=1$ uses the full band).

\paragraph{Radon transform.} The Radon transform $R_\phi(s)$ of $S$ integrates the spectrum along
lines at angle $\phi$. Along the direction of the null stripes the integration lines follow the
stripes, so the projection perpendicular to them oscillates with period $1/L$.
We compute it by rotating the spectrum and summing columns (verified to match
\texttt{skimage.transform.radon} to numerical precision and $\sim$10$\times$ faster),
for $\phi = 0,1,\dots,179^\circ$.

\paragraph{Cepstrum along each projection.} The 1-D Fourier transform of each projection is,
by the projection--slice theorem, a radial slice of the 2-D cepstrum
$C = \mathcal{F}^{-1}\{\log|B|\}$ (verified numerically in the test suite). For a box blur of
length $L$, $\log|\operatorname{sinc}|$ has period $1/L$ and its cepstrum has a strong
\emph{negative} peak at lag $L$. The (direction, lag) map
$P(\phi, q)$ therefore shows a single deep negative dip at $(\theta, L)$.

\paragraph{Angle, then length.} The Radon angle $\theta_R$ is the direction of maximum variance
of the projections of $S$ (DC masked, short rim chords ignored): integrating along the null
stripes keeps them separate, so that projection oscillates the most. The length is the lag of the
deepest negative dip of $P(\theta, q)$ for $q\in[q_{\min}, q_{\max}]$, refined by parabolic
interpolation. Our estimator searches $P$ jointly over $q$ and the directions within
$\pm3^\circ$ of $\theta_R$, then refines the angle at $0.25^\circ$ resolution, so a slightly
misplaced variance maximum is corrected by the cepstral evidence.

Phase~0 had suggested replacing the Radon angle by an unrestricted joint search over all
directions, on the grounds that strong strokes compete with the blur stripes in the variance. The
full evaluation showed the opposite: the unrestricted search is fooled by isolated deep dips in
the wrong direction more often than the variance rule is fooled by strokes. On the development
split the windowed refinement beat both alternatives (90.7\% success, vs.\ 89.5\% for the textbook
two-stage rule and 85.5\% for the unrestricted search; mean angle error $1.3^\circ$ vs.\
$1.4^\circ$ and $5.5^\circ$). This was noticed when the first test-split sweep ranked the
two-stage rule above the unrestricted search; the choice itself was then made on development data
only, and both alternatives are reported as reference rows.

\subsection{The text-aware stage}
\label{sec:textaware}
The proposal planned a text-contrast pre-processing step (local contrast normalisation and
edge enhancement) in front of the estimator. We implemented and evaluated nine such operators
(CLAHE, unsharp masking, Gaussian high-pass, Laplacian, difference of Gaussians, gradient
magnitude, local contrast normalisation, homomorphic filtering, and the proposed
LCN+Laplacian combination) on the identical sweep (Section~\ref{sec:preproc}). Contrast
normalisation was neutral, the derivative operators --- including the proposed combination ---
collapsed under noise because they raise exactly the high-frequency band where noise, JPEG residue
and resampling artefacts live and where the blur nulls are weakest, and only a mild Gaussian
high-pass helped on synthetic data. On real plates none of the contrast operators moved OCR
accuracy (Section~\ref{sec:design}).

The analysis pointed the other way. Plates and printed text in real photographs are small,
resampled and compressed regions, so the upper part of their spectrum carries little genuine
signal. The final text-aware stage is therefore a \emph{band-limited multi-band search}: the joint
search is run three times with band limits $c\in\{1, 0.5, 0.3\}$ and the estimate whose dip is
most \emph{directional} (largest $z_\angle$, below) is kept. Long blurs, whose nulls are densely
spaced at low frequencies, are then found in the clean low band; short blurs, whose first null
lies beyond $0.3$ Nyquist, are still found in the full band. The configuration was chosen on
development data only (Section~\ref{sec:design}).

\subsection{Confidence gate}
\label{sec:gate}
Deconvolution with a wrong kernel does not merely fail to sharpen: it amplifies the wrong
frequencies and leaves ringing and ghost strokes that make the plate \emph{less} legible. The gate
turns that silent failure into an explicit refusal (``not confident --- image returned
unchanged''). It combines three tests, each computed from quantities the estimator already has:
\begin{itemize}
  \item \textbf{Dip depth} (\emph{confidence}): depth of the winning dip against the background of
        the whole (direction, lag) map, in robust standard deviations (median/MAD).
  \item \textbf{Directional contrast} $z_\angle$: depth of the winning dip against the
        \emph{other directions at the same lag}. A linear-motion dip exists in one direction
        only, whereas isotropic structure --- defocus, soft or band-limited content, noise
        texture --- is deep in every direction and scores near zero.
  \item \textbf{Split-half agreement}: a uniform blur is the same everywhere, so the estimate is
        repeated on the left, right, top and bottom halves of the region; the number of halves
        (0--4) that reproduce the angle (within $10^\circ$) and length (within
        $\max(2\,\text{px}, 15\%)$) is counted.
\end{itemize}
A threshold of 0 switches a test off. The thresholds are calibrated on development data
(Section~\ref{sec:gateexp}); the resulting gate uses dip depth and split-half agreement.

\subsection{Non-blind restoration}
Richardson--Lucy (RL) is the primary restorer. Real images are not periodic and content outside
the frame is unknown, so we use the masked (unknown-boundary) RL variant: the estimate lives on a
canvas padded by the kernel radius and the data term is evaluated only on the observed pixels,
which removes the boundary ringing that periodic RL produces. For the synthetic sweeps, iterations
are stopped by the discrepancy principle (residual variance reaching the noise variance estimated
from the finest diagonal Haar detail band with the Donoho--Johnstone MAD estimator), because the best iteration count
depends strongly on noise (Section~\ref{sec:truepsf}). On real plates we use a fixed 30 iterations
as in the proposal. Wiener deconvolution with an edge-tapered canvas is implemented as the
comparison method.

\subsection{Evaluation-only OCR}
Legibility is measured with an off-the-shelf recogniser (RapidOCR, recognition model only) applied
to the plate after perspective rectification with the dataset's corner annotations. It is a
measuring instrument and never part of the restoration. Chinese plates start with a province
glyph; we score the six alphanumeric characters (character accuracy $=1-$ edit distance$/6$).
On sharp CCPD plates the instrument reads 100\% of the characters correctly.
